\section{Proofs for conditional recourse}\label{app:rec}

This appendix proves the results of \cref{sec:rec}. \Cref{app:rec:tools}
lists the external tools. \Cref{app:rec:search,app:rec:certificates,app:rec:convex}
prove the common machinery, \cref{app:rec:main} proves the ambient-noise
theorems, \cref{app:rec:core-only} proves the core-only theorem, and
\cref{app:rec:examples} verifies the examples.

\subsection{Conventions and tools}\label{app:rec:tools}

Norms are Euclidean unless marked otherwise. For a symmetric matrix the
operator norm is at most the largest absolute row sum. Hence, with $K_2,K_3$
as in \eqref{eq:rec:K}, every principal submatrix and every off-diagonal
block of $\nabla^2F_0(x)$ has operator norm at most $K_2$, and
\begin{equation}\label{eq:app:rec:hessvar}
 \norm{\nabla^2F_0(x)-\nabla^2F_0(y)}\le K_3\norm{x-y}_\infty,\qquad
 \abs{D^3F_0(x)[p,p',p'']}\le K_3\norm p\norm{p'}\norm{p''}
\end{equation}
for $x,y\in[0,1]^n$ and vectors $p,p',p''$. For the first bound, the
$(i,j)$ entry of the difference is at most
$\sum_l\sup\abs{\partial_{ijl}F_0}\abs{x_l-y_l}$, so every row sum is at
most $K_3\norm{x-y}_\infty$. For the second, the symmetric matrix
$(\sum_l\partial_{ijl}F_0\,p''_l)_{ij}$ has row sums at most
$K_3\norm{p''}_\infty\le K_3\norm{p''}$. The same row-sum argument gives
$\abs{\partial_iF_0(x)-\partial_iF_0(y)}\le K_2\norm{x-y}_\infty$. Monomial
coefficient sums give rational $K_2,K_3,G$ of polynomial bit length for
fixed $d$.

We use the following results. Items \ref{it:app:rec:growth}--\ref{it:app:rec:rare}
are proved in \cref{sec:count} and its appendix; we restate the forms used.

\begin{enumerate}[label=(T\arabic*),leftmargin=*]
\item\label{it:app:rec:local} \emph{Corrected corners and local counts}
(\cref{lem:count:rounding,thm:count:local}). If $f$ has upper coordinate
curvature at most $L$ on a box and $C$ is a cell with side lengths
$\eta_i$, then $\min_{\rm corners}F_c\le F_c(x)+\frac18\sum_iL\eta_i^2$ for
every $x\in C$ and every $c$. For the uniform grid with $m_i$ parts of width
$h_i$ per coordinate and independent coefficients with interval
concentrations $\pi_i$, the expected number of nodes $v$ with
$F_c(v)\le F_c(v\pm h_ie_i)+\delta$ for all interior coordinates is at most
$\prod_i[2+(m_i-1)\min\{1,\pi_i(Lh_i+2\delta/h_i)\}]$, also conditionally on
any random element independent of $c$ on which $f$ may depend.
\item\label{it:app:rec:growth} \emph{Growth tail}
(\cref{thm:count:growth-tail}). If $Z\subset\R^p$ is compact with coordinate
widths $w_i$, $f:Z\to\R$ is lower semicontinuous, and $c$ is uniform on
$[-\sigma,\sigma]^p$, then the point-growth modulus of $f+c^{\mathsf T}y$
satisfies $\Pr\{g_*<\varepsilon\}\le\varepsilon\sum_iw_i/\sigma$.
\item\label{it:app:rec:finite} \emph{Finite-law tails}
(\cref{lem:count:finite-tails}). Let $\mathcal D\subset\R^{p+m}$ be compact
and described by at most $s$ polynomial sign conditions of degree at most
$d$ (bounded integer coordinates by disjunctions), let $F$ have degree at
most $d$, and perturb only the first $p$ coordinates. For the reduced
objective $y\mapsto\min\{F(y,y'):(y,y')\in\mathcal D\}$ on the projection of
$\mathcal D$, with coefficients independent with law $U_{\sigma,M}$,
$\Pr\{g_*<\varepsilon\}\le S\varepsilon/\sigma+2pC_{\rm tail}/M$, where $S$ is
the sum of the projected widths and $C_{\rm tail}$ is base-computable with
$\log C_{\rm tail}\le\poly_d(I)$ when $\log(s+1)\le\poly(I)$. If $m=0$ and
$\mathcal D=[0,1]^n$, the probability that $g_*>0$ and the minimizer has a
coordinate at a bound with $\abs{\partial_iF_\gamma}\le\tau$ is at most
$N_{\rm act}(\tau/\sigma+1/M)$, $N_{\rm act}=n3^n\max\{1,d-1\}^n$.
\item\label{it:app:rec:fallback} \emph{Exact fallback}
(\cref{thm:count:fallback}). For a fixed-degree polynomial on an explicit
compact mixed semialgebraic domain with base length $I$, there are a
base-computable $B=2^{\poly_d(I)}$ and a fixed exponent such that, for all
rational linear coefficients of bit length at most $b$, a deterministic
algorithm returns the optimal value and the lexicographically least global
minimizer, each coordinate in root representation, within $B\poly_d(I+b)$
bit operations, and refinements to $q$ bits within $B\poly_d(I+b+q)$.
\item\label{it:app:rec:rare} \emph{Rare fallback}
(\cref{lem:count:rare-fallback}(a)). If the fallback costs at most
$B\,P(L_{\max})$ on every draw and runs with probability at most $1/B$, it
contributes at most $P(L_{\max})$ to the expected work. The parameters are
chosen in the order $B$, thresholds, stopping depth $J$, grid size $M$.
\item\label{it:app:rec:qe} \emph{Fixed-block quantifier elimination}
\citep[Theorem~1.1]{renegar1992-on-the-computational-complexity-and}. A formula with $\omega$ quantifier blocks
of sizes $n_1,\ldots,n_\omega$, $\ell$ free variables, $s$ atoms, and degree
at most $d_0$, with arbitrary real coefficients, is equivalent to a
quantifier-free formula in disjunctive normal form whose numbers of
disjuncts and of atoms per disjunct, and whose atom degrees, are at most
$(sd_0)^{2^{c\omega}\ell\prod_bn_b}$ for an absolute constant $c$. If the
coefficients are integers of bit length at most $\tau$, the output
coefficients have bit length at most $(\tau+1)(sd_0)^{2^{c\omega}\ell\prod_bn_b}$.
\item\label{it:app:rec:tube} \emph{Tubes around algebraic sets}
\citep[Theorem~1.1]{basu2023-hausdorff-approximations-and-volume-of}. If $Z\subseteq\R^q$ is a real algebraic
set defined by polynomials of degree at most $D$ with $\dim Z\le m$, and $X$
is uniform in a Euclidean ball of radius $T$, then for every $\epsilon>0$
\[
 \Pr\{\dist(X,Z)\le\epsilon\}
 \le4\Bigl(\frac{4qD\epsilon}T\Bigr)^{q-m}\Bigl(1+\frac{(4D+1)\epsilon}T\Bigr)^m .
\]
\item\label{it:app:rec:gls} \emph{Convex values and patch evaluation}
(\cref{lem:sp:gls}, based on weak convex optimization
\citep[Corollary~4.2.7]{grotschel1988-geometric-algorithms-and-combinatorial-optimization}). For a convex
polynomial $f$ of fixed degree on a rational box $P$ with nonempty interior
relative to its free coordinates, and a rational $\eta\in(0,1]$, one computes
a rational $y\in P$ and rationals $a\le\min_Pf\le f(y)\le U\le a+\eta$ in
time polynomial in the length of the data and $\log(1/\eta)$; here the
oracle of that lemma is exact rational evaluation of $f$ and $\nabla f$, and
the repair map is coordinatewise clipping to $P$, which is nonexpansive. If
$\nabla^2f\succeq gI$ on $P$ with rational $g>0$, the unique minimizer $y^*$
satisfies $\norm{y-y^*}\le\sqrt{2\eta/g}$.
\item\label{it:app:rec:bezout} \emph{Nonsingular zeros} (\cref{lem:sp:bezout}).
A system of $q$ polynomial equations of degree at most $D$ in $q$ complex
variables has at most $D^q$ zeros with nonsingular Jacobian.
\item\label{it:app:rec:faces} \emph{Quadratic faces} (\cref{lem:sp:faces}).
For a rational quadratic on a mixed box with $n_c$ continuous coordinates
and $N_Z$ integer labels, solving the stationary system of every face whose
free Hessian block is positive definite yields a rational global minimizer,
in $3^{n_c}N_Z\poly(I+b)$ bit operations.
\end{enumerate}

\subsection{Conditional values and the core search}\label{app:rec:search}

\begin{proof}[Proof of \cref{lem:rec:value}]
(a) The term $\gamma_C^{\mathsf T}v$ does not depend on $z$, so
$V(v)-\gamma_C^{\mathsf T}v=\min_{z\in Y}[F_0(v,z)+\gamma_R^{\mathsf T}z]$.

(b) Fix $z\in Y$. Along the segment $t\mapsto(v+te_i,z)$, which stays in
$[0,1]^k\times\conv Y$, the polynomial $\phi_z(t)=F_\gamma(v+te_i,z)$ has
$\phi_z''=\partial_{ii}F_0\le L$, so $\phi_z(t)-Lt^2/2$ is concave. The
function $V(v+te_i)-Lt^2/2=\min_{z\in Y}[\phi_z(t)-Lt^2/2]$ is a pointwise
minimum of concave functions over an index set that does not depend on $t$,
hence concave.

(c) Let $z\in Y'$ attain $V_{Y'}(w)$; it exists because $Y'$ is compact. The
segment from $(w,z)$ to $(v,z)$ lies in $[0,1]^k\times\conv Y$, so the mean
value theorem and H\"older's inequality give
$V_{Y'}(v)\le F_\gamma(v,z)\le F_\gamma(w,z)+G\norm{v-w}_\infty
=V_{Y'}(w)+G\norm{v-w}_\infty$. Exchange $v$ and $w$.
\end{proof}

By \cref{lem:rec:value}(a),(b), the function $\bar V=V-\gamma_C^{\mathsf T}v$
has upper coordinate curvature at most $L$ on $[0,1]^k$ and does not depend
on $\gamma_C$. Applied to $f=\bar V$ and $c=\gamma_C$, the first part of
\ref{it:app:rec:local} gives \eqref{eq:rec:corner}: every point $v$ of a
level-$j$ cell has a corner $c$ with $V(c)\le V(v)+E_j$.

\begin{proof}[Proof of \cref{prop:rec:search}]
Write $\ell(v),u(v)$ for the oracle's answers at a queried corner, with
$\ell=u=V$ in exact mode, and put $\varepsilon'_j=0$ in exact mode and
$\varepsilon'_j=E_j$ in certified mode. Then
$\ell(v)\le V(v)\le u(v)\le\ell(v)+\varepsilon'_j$ and
$\varepsilon_j=E_j+\varepsilon'_j$. Level $0$ generates the cell
$[0,1]^k$, queries its corners, and applies \eqref{eq:rec:retain}; level
$j\ge1$ generates the children of the cells retained at level $j-1$.

(a) Let $(a,z^*)$ be a global minimizer; then $V(a)=f^*$. Suppose a cell
generated at level $j$ contains $a$; at level $0$ this is $[0,1]^k$. By
\eqref{eq:rec:corner} some corner $c$ has $V(c)\le f^*+E_j$, so
$\min_{\rm corners}\ell-E_j\le V(c)-E_j\le f^*$. Every upper value is the
value of a feasible point, so $U_j\ge f^*$, and the cell is retained. A
child of it contains $a$ and is generated at level $j+1$. Induction proves
(a).

(b) With $c$ as in (a), $U_j\le u(c)\le V(c)+\varepsilon'_j\le
f^*+\varepsilon_j$. The cell that generated the corner $c_j$ satisfies
$\min_{\rm corners}\ell-E_j\le\ell(c_j)\le u(c_j)=U_j$, so it is retained.

(c) A retained cell has a corner $v$ with $\ell(v)-E_j\le U_j$, so
$V(v)\le\ell(v)+\varepsilon'_j\le U_j+E_j+\varepsilon'_j\le f^*+2\varepsilon_j$.

\emph{Count.} Fix $j\le J$ and $h=h_j$, and let $N_j$ be the number of
points $v$ of the deterministic grid $(h\Z)^k\cap[0,1]^k$ with
$V(v)\le f^*+2\varepsilon_j$. Such a point satisfies
$V(v)\le V(v\pm he_i)+2\varepsilon_j$ for every coordinate with $0<v_i<1$,
because its grid neighbors are feasible. Apply the second part of
\ref{it:app:rec:local} to $f=\bar V$, $c=\gamma_C$, $m_i=2^j$, $h_i=h$, and
$\delta=2\varepsilon_j$, conditionally on $\gamma_R$, which is independent of
$\gamma_C$. The relevant interval length is $Lh+4\varepsilon_j/h$, which equals
$\lambda h$ with $\lambda=L(1+k/2)$ in exact mode and $\lambda=L(1+k)$ in
certified mode, and \eqref{eq:model:interval} gives
$\pi_i(\lambda h)\le\lambda h/(2\sigma)+1/M$. Since $h^{-1}=2^j\le2^J\le M$,
\[
 \E[N_j\mid\gamma_R]\le\prod_{i=1}^k\Bigl[2+(h^{-1}-1)\Bigl(\frac{\lambda h}{2\sigma}
 +\frac1M\Bigr)\Bigr]\le\Bigl[3+\frac\lambda{2\sigma}\Bigr]^k,
\]
which is $Q_{\rm ex}$ or $Q_{\rm ap}$.

\emph{Oracle calls.} By (c), each retained cell at level $j$ has a corner
counted by $N_j$, and a grid point is a corner of at most $2^k$ cells of
level $j$; so at most $2^kN_j$ cells are retained. Level $j\ge1$ generates
at most $2^k\cdot2^kN_{j-1}$ cells and queries at most $8^kN_{j-1}$ corners;
level $0$ queries $2^k$ corners. Taking expectations gives at most
$2^k+8^kJQ$ calls through level $J$.
\end{proof}

The count compares the adaptive search with a deterministic grid fixed before
sampling: retained cells are charged to grid points, and the event that a grid
point is near-optimal is stated through the true function $V$. The oracle's
answers may depend on all noise coefficients; they enter only through
(a)--(c), which hold for every admissible answer.

\subsection{Containment and closure}\label{app:rec:certificates}

\begin{proof}[Proof of \cref{lem:rec:exclusion}]
\emph{Cover.} If $z\in Y\setminus P$, some coordinate violates its patch
interval. If $z_i<\max\{\ell_i,z_{c,i}-r\}$, then $z_i<z_{c,i}-r$ because
$z_i\ge\ell_i$; hence $\ell_i\le z_i<z_{c,i}-r$, the lower slab $Y_i^-$ is
formed, and it contains $z$. The upper case is symmetric.

\emph{Containment.} Let $v\in\mathcal H$ and let $z$ minimize
$F_\gamma(v,\cdot)$ on $Y$. If $z\notin P$, then $z$ lies in some slab $S$.
Since $c,v\in\mathcal H$, $\norm{v-c}_\infty\le w_\infty(\mathcal H)$. By
\cref{lem:rec:value}(c) applied to $S$ and to $Y$,
\[
 V(v)=F_\gamma(v,z)\ge V_S(v)\ge V_S(c)-Gw_\infty(\mathcal H)
 \ge\lambda_{\rm out}-Gw_\infty(\mathcal H)>U+Gw_\infty(\mathcal H)
 \ge V(c)+Gw_\infty(\mathcal H)\ge V(v),
\]
a contradiction. If $(a,z^*)$ is a global minimizer, then $z^*$ minimizes
$F_\gamma(a,\cdot)$ on $Y$; when $a\in\mathcal H$ it follows that $z^*\in P$.
\end{proof}

\begin{proof}[Proof of \cref{lem:rec:closure}]
(a) Let $x^*$ be a global minimizer; it lies in $Q$. If $\partial_iF_\gamma>0$
on $Q$, the original lower bound $\ell_i$ lies in the $i$th interval $Q_i$,
and $x^*_i>\ell_i$, then the segment from $x^*$ to
$x^*-(x^*_i-\ell_i)e_i$ lies in $Q\subseteq X$ and $F_\gamma$ strictly
decreases along it, which is impossible. The midpoint test is sufficient
because $\abs{\partial_iF_\gamma(x)-\partial_iF_\gamma(m)}\le K_2r_Q$ on $Q$
(\cref{app:rec:tools}). For a quadratic, $\partial_iF_\gamma$ is affine and
its extreme values on $Q$ are attained at vertices chosen coordinatewise.
Each fixing is justified with respect to the same box $Q$, so all of them
hold simultaneously for every global minimizer.

(b) Let $\mathcal U$ be the set of unfixed coordinates. For $x\in Q'$,
\eqref{eq:app:rec:hessvar} gives
$\nabla^2_{\mathcal U\mathcal U}F_0(x)\succeq\nabla^2_{\mathcal U\mathcal U}
F_0(m')-K_3r_{Q'}I\succeq g_0I$. Thus $F_\gamma$ is $g_0$-strongly convex on
the box $Q'$ and has a unique minimizer there. By (a), $Q'\subseteq X$
contains every global minimizer, so $\min_{Q'}F_\gamma=f^*$, and the unique
minimizer on $Q'$ is the unique global minimizer.

(c) As in (b), $\min_{Q'}F_\gamma=f^*$, and the restriction is a convex
quadratic, so every minimizer on $Q'$ attains $f^*$.
\end{proof}

\subsection{Certified convex recourse}\label{app:rec:convex}

\begin{proof}[Proof of \cref{lem:rec:convex-oracle}]
Fix a query $(v,Y',\gamma,\eta)$ and let $f=F_\gamma(v,\cdot)$ on $Y'$, after
substituting the coordinates whose interval is a single point; if none
remains, return the value. Otherwise $f$ is a convex polynomial on the box
$B=Y'$ with explicit rational coefficients of polynomial length, and monomial
bounds give a rational
$K_f\ge\max\{1,\max_B\norm{\nabla^2f}\max(1,\diam B)^2\}$ of polynomial
length. Let $\delta=\min\{\eta/2,\eta^2/(8K_f),1\}$.

\emph{A feasible near-minimizer.} By \ref{it:app:rec:gls}, with accuracy
$\delta$, one computes a rational $y\in B$ with $f(y)\le\min_Bf+\delta$. That
lemma uses weak convex optimization over a capped epigraph and repairs the
near-feasibility of its output by clipping; it uses convexity of $f$ only.

\emph{The tangent certificate.} Put $u=f(y)$ and
$\ell=f(y)+\min_{z\in B}\nabla f(y)^{\mathsf T}(z-y)$, an explicit choice of one
endpoint per coordinate. By convexity $f(z)\ge f(y)+\nabla f(y)^{\mathsf T}(z-y)
\ge\ell$ on $B$, so $\ell\le\min_Bf$. Let $\bar z$ attain the linear minimum
and $\Delta=u-\ell=-\nabla f(y)^{\mathsf T}(\bar z-y)\ge0$. For $t\in(0,1]$ the
point $y+t(\bar z-y)$ lies in $B$, and Taylor's theorem gives
$\min_Bf\le f(y)-t\Delta+K_ft^2/2$. Since $f(y)\le\min_Bf+\delta$,
$\Delta\le\delta/t+K_ft/2$. Choose $t=\min\{1,\eta/(2K_f)\}$. If
$\eta\le2K_f$, the two terms are each at most $\eta/4$; otherwise $t=1$ and
$\Delta\le\eta/2+\eta/4$. In both cases $u-\ell\le\eta$. All quantities have
bit length polynomial in the query length and the length of $\eta$, so the
bit work is polynomial. A verifier recomputes $f(y)$, $\nabla f(y)$, and the
endpoint minimum, and needs neither the optimization trace nor a modulus of
strong convexity.
\end{proof}

\subsection{Proofs of the ambient-noise theorems}\label{app:rec:main}

We prove \cref{thm:rec:qp,thm:rec:poly} together. In the quadratic case
$F_0=\frac12x^{\mathsf T}Ax+a^{\mathsf T}x+c_0$, the oracle runs in exact mode,
$d=2$, and $K_3$ is not used.

\paragraph{Base choices.}
Compute $K_2,K_3$ from \eqref{eq:rec:K} and
$G=\max\{1,\sum_{i\in C}(\sup_{[0,1]^n}\abs{\partial_iF_0}+\sigma)\}$, which
satisfies the hypothesis of \cref{lem:rec:value}(c) for every noise vector
in $[-\sigma,\sigma]^n$. Let $B=\max\{2,3^n\}$ in the quadratic case
(\ref{it:app:rec:faces} with $N_Z=1$) and the factor of \ref{it:app:rec:fallback} in the
polynomial case, enlarged so that it also bounds $q$-bit evaluation. Let
$N_{\rm act}$ and $C_{\rm tail}$ be as in \ref{it:app:rec:finite} for
$[0,1]^n$, and choose
\begin{equation}\label{eq:app:rec:params}
\begin{gathered}
 \rho=\frac1{4B},\qquad g_0=\frac{\rho\sigma}{2n},\qquad
 \tau=\frac{\rho\sigma}{2N_{\rm act}},\qquad
 r=\min\Bigl\{\frac18,\frac\tau{16K_2},\frac{g_0}{4K_3}\Bigr\},\qquad
 A_0=2+\frac{kL+8}{g_0},\\
 J=\min\Bigl\{j\ge0:\ 2^{-j}\le\min\Bigl\{\frac r{4A_0},\
 \frac{g_0r^2}{16GA_0}\Bigr\}\Bigr\},\qquad
 M=\text{least power of two}\ \ge\max\Bigl\{2,2^J,\frac{4nC_{\rm tail}}\rho,
 \frac{2N_{\rm act}}\rho\Bigr\}.
\end{gathered}
\end{equation}
In the quadratic case omit the term $g_0/(4K_3)$. All these numbers are
computed from the base instance, have bit length polynomial in $I$, and
$J,\log_2M\le\poly_d(I)$.

\paragraph{The algorithm.}
For $j=0,1,\ldots,J$:
\begin{enumerate}[label=(\arabic*)]
\item Run level $j$ of the core search: exact mode in the quadratic case,
certified mode with accuracy $E_j$ otherwise. If $k=0$, make one oracle call
with accuracy $4^{-j}$, let $U_j$ be its upper value, and let
$\mathcal H_j$ be the empty core point.
\item Form the patch $P$ of radius $r$ about $z_j$ and its slabs. For each
slab obtain its exact value (quadratic case) or a certified lower bound with
accuracy $\eta_{\rm out}=g_0r^2/16$; let $\lambda_{\rm out}$ be the least of
these. If \eqref{eq:rec:exclusion} fails with $U=U_j$ and
$\mathcal H=\mathcal H_j$, go to the next level.
\item Apply the sign tests of \cref{lem:rec:closure}(a) on
$Q=\mathcal H_j\times P$ and obtain $Q'$. If no coordinate is unfixed,
return the fixed point and its value. In the quadratic case test whether
the Hessian block of the unfixed coordinates is positive semidefinite; if so,
solve the convex quadratic program on $Q'$ exactly
\citep{kozlov1980-the-polynomial-solvability-of-convex} and return its solution. In the
polynomial case test \cref{lem:rec:closure}(b) with $g_0$; if it passes,
return $(S,z_S,Q',g_0)$. Otherwise go to the next level.
\end{enumerate}
If no level returns, run the fallback on the same draw:
\ref{it:app:rec:faces} in the quadratic case and \ref{it:app:rec:fallback}
otherwise.

\paragraph{Correctness on every draw.}
By \cref{prop:rec:search}(a),(b), $\mathcal H_j$ contains the core of every
global minimizer and the corner $c_j$, and $U_j\ge V(c_j)$. Each certified
lower bound is at most the true slab value. Hence \cref{lem:rec:exclusion}
shows that every global minimizer lies in $\mathcal H_j\times P$ when step~(2)
passes, and \cref{lem:rec:closure} shows that the outputs of step~(3) are
correct. The fallback is exact. No step assumes growth, uniqueness, or any
probabilistic event.

\begin{lemma}[Stopping]\label{lem:app:rec:stopping}
Suppose $F_\gamma$ has a unique global minimizer $x^*=(a,z^*)$ with point
growth at least $g_0$, and every coordinate of $x^*$ at a bound has
$\abs{\partial_iF_\gamma(x^*)}>\tau$. Then the algorithm returns at a level
$j\le J$.
\end{lemma}

\begin{proof}
Assume no earlier level returned; we show that level $J$ returns. Put
$h=2^{-J}$ and $E=E_J$.

\emph{Hull and incumbent.} Every retained cell has a corner $v$ with
$V(v)\le f^*+4E$ (\cref{prop:rec:search}(c); $2E$ in exact mode), and
$V(v)-f^*\ge g_0\norm{v-a}^2$. Hence
$\norm{v-a}\le h\sqrt{kL/(2g_0)}$, and every point of the cell is within
$h$ of $v$ in each coordinate. Thus $\mathcal H_J\subseteq a+[-A_0h,A_0h]^k$,
because $1+\sqrt y\le2+y$. The incumbent value satisfies
$\varepsilon_{\rm inc}=U_J-f^*\le2E$ for $k\ge1$ and
$\varepsilon_{\rm inc}\le h^2$ for $k=0$, and growth gives
$\norm{(c_J,z_J)-x^*}^2\le\varepsilon_{\rm inc}/g_0$. Since
$A_0\ge\max\{2,kL/g_0,8/g_0\}$, this yields $\norm{(c_J,z_J)-x^*}\le A_0h\le r/4$
and $\varepsilon_{\rm inc}\le g_0r^2/16$; for $k\ge1$ the last inequality is
$kLh^2/4\le kLr^2/(64A_0^2)\le g_0r^2/16$, and for $k=0$ it is
$h^2\le r^2/(16A_0^2)\le g_0r^2/256$.

\emph{Containment.} A point $z$ of a slab differs from $z_J$ by at least $r$
in one coordinate, hence from $z^*$ by at least $3r/4$. Growth gives
$F_\gamma(c_J,z)-f^*\ge9g_0r^2/16$, so
$\lambda_{\rm out}\ge f^*+9g_0r^2/16-\eta_{\rm out}=f^*+g_0r^2/2$ (or
$f^*+9g_0r^2/16$ in exact mode). Therefore
$\lambda_{\rm out}-U_J\ge7g_0r^2/16$, while
$2Gw_\infty(\mathcal H_J)\le4GA_0h\le g_0r^2/4$. Step~(2) passes.

\emph{Sign tests.} Every point of $Q=\mathcal H_J\times P$ is within $5r/4$
of $x^*$ in the maximum norm, and $r_Q\le r$. Let $x^*_i$ be at a bound, say
the lower one, with $\partial_iF_\gamma(x^*)>\tau$; that bound lies in
$Q_i$ because $x^*\in Q$. The midpoint test value is at least
$\tau-\frac54K_2r-K_2r\ge\tau-\frac9{64}\tau>0$, so the coordinate is fixed;
in the quadratic case the exact test passes because $\partial_iF_\gamma\ge
\tau-\frac54K_2r>0$ on $Q$. A coordinate strictly inside its bounds at $x^*$
has zero derivative at the point $x^*\in Q$ and cannot pass a strict sign
test. Hence $Q'$ fixes exactly the coordinates that are at a bound at $x^*$,
and the unfixed set $\mathcal U$ consists of the coordinates strictly inside.

\emph{Hessian test.} For $p$ supported on $\mathcal U$ and small $t$,
$F_\gamma(x^*+tp)-f^*\ge g_0t^2\norm p^2$ and $\nabla_{\mathcal U}F_\gamma(x^*)=0$,
so $\nabla^2_{\mathcal U\mathcal U}F_0(x^*)\succeq2g_0I$. In the quadratic case
the semidefinite test passes. Otherwise $x^*\in Q'$ gives
$\norm{m'-x^*}_\infty\le r_{Q'}\le r$, and \eqref{eq:app:rec:hessvar} gives
$\nabla^2_{\mathcal U\mathcal U}F_0(m')-(K_3r_{Q'}+g_0)I\succeq
(g_0-2K_3r)I\succeq\frac{g_0}2I$ because $r\le g_0/(4K_3)$.
\end{proof}

\begin{proof}[Proof of \cref{thm:rec:qp,thm:rec:poly}]
Correctness was shown above. Let $\mathcal B$ be the event that the
hypothesis of \cref{lem:app:rec:stopping} fails. By \ref{it:app:rec:growth}
and \ref{it:app:rec:finite} applied to $F_0$ on $[0,1]^n$ with ambient
noise,
\[
 \Pr\{g_*<g_0\}\le\frac{ng_0}\sigma+\frac{2nC_{\rm tail}}M\le\frac\rho2+\frac\rho2,
 \qquad
 \Pr\{g_*\ge g_0,\ \text{a small active gradient}\}
 \le N_{\rm act}\Bigl(\frac\tau\sigma+\frac1M\Bigr)\le\rho,
\]
so $\Pr(\mathcal B)\le2\rho=1/(2B)$, and the fallback runs only on
$\mathcal B$.

\emph{Work.} Since $M\ge2^J$, \cref{prop:rec:search} bounds the expected
number of core oracle calls by $2^k+8^kJQ$. Every query has length
$\poly_d(I)$: core corners are dyadic with at most $J$ bits per coordinate,
sampled coefficients have $O(I+\log M)$ bits, the accuracies $E_j$ and
$\eta_{\rm out}$ have polynomial length, and the oracle's outputs, hence the
patch endpoints $z_{j,i}\pm r$, have polynomial length. Each oracle call
therefore costs $\poly_d(I)$, with an exponent independent of $k$. Each level
adds at most $2n_R$ slab calls and $\poly_d(I)$ work for the sign and Hessian
tests and the final convex solve. The fallback costs $B\poly_d(I+\log M)$ and
runs with probability at most $1/(2B)$ (\ref{it:app:rec:rare}). Summing,
the expected work is $[2^k+8^kJQ+(J+1)n_R]\poly_d(I)\le8^kQ\poly_d(I)$, with
$Q=Q_{\rm ex}$ or $Q_{\rm ap}$. The global proof record consists of the oracle
answers with their checkable certificates (exact values, or tangent bounds
\eqref{eq:rec:tangent}), the retained lists, the slab bounds, and the test
data; its expected size obeys the same bound.

\emph{Outputs.} In the quadratic case the returned point is rational of
polynomial length on every draw, by \citet{kozlov1980-the-polynomial-solvability-of-convex} on the
ordinary branch and \ref{it:app:rec:faces} on the fallback branch. In the
polynomial case the descriptor $(S,z_S,Q',g_0)$ consists of dyadic hull
endpoints, patch endpoints of polynomial length, and the base number $g_0$;
\ref{it:app:rec:gls}, run with $\eta=\min\{2^{-q},(g_0/2)2^{-2q}\}$, evaluates it in $\poly_d(I+q)$. The fallback output
has $q$-bit evaluation cost $B\poly_d(I+\log M+q)$ and occurs with
probability at most $1/(2B)$, which gives the expected evaluation bound.
\end{proof}

\begin{proof}[Proof of \cref{cor:rec:forest}]
Given a query $(v,\prod_{i\in R}[\ell'_i,u'_i],\gamma)$, substitute $v$ and
every residual coordinate whose interval is a point. For each remaining
coordinate write $z_i=\ell'_i+(u'_i-\ell'_i)y_i$ with $y_i\in[0,1]$. The
result is $y^{\mathsf T}\tilde Ay+\tilde c^{\mathsf T}y+\text{const}$ with
$\tilde A_{ij}=\frac12A_{ij}(u'_i-\ell'_i)(u'_j-\ell'_j)$, so its interaction
graph is a subgraph of the residual graph, a forest. The algorithm of
\citet[Theorem~1 and Lemma~19]{pia2026-treewidth-and-the-complexity-of} returns a rational
optimal $y$ of polynomial bit length using $O(n^2)$ arithmetic operations on
numbers of polynomial length, hence in polynomial bit time with an absolute
exponent. Mapping back and evaluating gives the oracle's answer. Diagonal
entries of $A_{RR}$ of either sign are allowed.
\end{proof}

\begin{proof}[Proof of \cref{prop:rec:rank}]
The Hessian blocks are $\partial_{tt}F=\lambda$,
$\nabla_z\partial_tF=T\nabla P(z)$, and $\nabla^2_{zz}F=Tt\nabla^2P(z)$ with
$\nabla^2P=12\diag(z_i^2)+12(\sum_iz_i)^2\mathbf1\mathbf1^{\mathsf T}\succeq0$.
So the residual problem is convex for every $t\in[0,1]$, also after linear
perturbations and box restrictions.

Suppose $F+\frac12(t,z)^{\mathsf T}K(t,z)$ is convex on the box. Its Hessian
$\nabla^2F+K$ is positive semidefinite in the interior and, by continuity, on
the closed box, in particular at $t=0$, where
$\nabla^2F(0,z)=\bigl(\begin{smallmatrix}\lambda&T\nabla P(z)^{\mathsf T}\\
T\nabla P(z)&0\end{smallmatrix}\bigr)$. Let $u\in\ker K_{zz}$ and $w=(0,u)$.
Then $w^{\mathsf T}Kw=0$, so $Kw=0$, and
$w^{\mathsf T}(\nabla^2F(0,z)+K)w=0$, so $(\nabla^2F(0,z)+K)w=0$ and
$\nabla^2F(0,z)w=0$. Its first component gives $T\nabla P(z)^{\mathsf T}u=0$
for all $z\in[0,1]^m$. Since $\nabla P(e_i)=4(e_i+\mathbf1)$ and the matrix
$I+\mathbf1\mathbf1^{\mathsf T}$ is invertible, $u=0$. Thus the positive
semidefinite matrix $K_{zz}$ is nonsingular, and $\rank K\ge m$.

For the remark after the proposition: on points with $t>0$ and all $z_i>0$,
$\nabla^2P\succ0$, and Euler's identities for the quartic form $P$ give
$\nabla^2P(z)z=3\nabla P(z)$ and $\nabla P(z)^{\mathsf T}z=4P(z)$, hence
$\nabla P^{\mathsf T}(\nabla^2P)^{-1}\nabla P=\frac43P$. Adding $Ct\log t$ makes
the Schur complement of the residual block equal to
$\lambda+C/t-\frac{4T}{3t}P(z)\ge\lambda$ when $C\ge\frac43T\max P$. Continuity
extends convexity to the closed box, where $t\log t$ is extended by $0$ at
$t=0$.
\end{proof}

\subsection{Core-only noise}\label{app:rec:core-only}

Throughout, the hypotheses of \cref{thm:rec:core-only} hold, $Y=[0,1]^{n_R}$,
$H=K_2/\mu$, and $H_V=K_2(1+H)$. The residual selector is
$s(v)=\argmin_{z\in Y}F_0(v,z)$, and $V_0(v)=F_0(v,s(v))$. Core-only noise
changes neither $s$ nor the residual active sets, and
$V_\gamma=V_0+\gamma^{\mathsf T}v$.

\begin{lemma}[Selector, value, and growth lift]\label{lem:app:rec:lift}
\leavevmode
\begin{enumerate}[label=(\alph*)]
\item The selector is unique and $\norm{s(v)-s(w)}\le H\norm{v-w}$.
\item $V_0$ is differentiable on $[0,1]^k$ (one-sided at the boundary), with
$\nabla V_0(v)=\nabla_vF_0(v,s(v))$, and $\nabla V_0$ is $H_V$-Lipschitz.
\item If $V_\gamma(v)-V_\gamma(a)\ge g_V\norm{v-a}^2$ on $[0,1]^k$ with
$g_V>0$, then $x^*=(a,s(a))$ is the unique global minimizer of $F_\gamma$ and
$F_\gamma(x)-F_\gamma(x^*)\ge g_F\norm{x-x^*}^2$ on $X$, where
$g_F=\min\{g_V/(1+2H^2),\mu/4\}$.
\end{enumerate}
\end{lemma}

\begin{proof}
(a) Uniqueness follows from strong convexity. Write $\nabla_z$ for the
residual gradient. The optimality conditions
$\nabla_zF_0(v,s(v))^{\mathsf T}(y-s(v))\ge0$ and
$\nabla_zF_0(w,s(w))^{\mathsf T}(y-s(w))\ge0$ for $y\in Y$, with $y=s(w)$ and
$y=s(v)$ respectively, add to
$[\nabla_zF_0(w,s(w))-\nabla_zF_0(v,s(v))]^{\mathsf T}(s(w)-s(v))\le0$. Split
the bracket at $(w,s(v))$. Strong monotonicity bounds the first part below by
$\mu\norm{s(w)-s(v)}^2$, and the cross block of the Hessian bounds the second
part by $K_2\norm{w-v}\norm{s(w)-s(v)}$.

(b) Taylor's theorem in $v$ at fixed residual points gives
$V_0(w)\le F_0(w,s(v))\le V_0(v)+\nabla_vF_0(v,s(v))^{\mathsf T}(w-v)
+\frac{K_2}2\norm{w-v}^2$ and
$V_0(w)=F_0(w,s(w))\ge V_0(v)+\nabla_vF_0(v,s(w))^{\mathsf T}(w-v)
-\frac{K_2}2\norm{w-v}^2$. By (a),
$\norm{\nabla_vF_0(v,s(w))-\nabla_vF_0(v,s(v))}\le K_2H\norm{w-v}$. Hence
$\abs{V_0(w)-V_0(v)-\nabla_vF_0(v,s(v))^{\mathsf T}(w-v)}\le
K_2(\frac12+H)\norm{w-v}^2$, which proves differentiability. The Lipschitz
bound follows by splitting at $(w,s(v))$ and using (a).

(c) For $x=(v,z)\in X$, strong convexity in $z$ and the optimality of
$s(v)$ on the box give $F_\gamma(v,z)-V_\gamma(v)\ge\frac\mu2\norm{z-s(v)}^2$,
so $F_\gamma(x)-f^*\ge g_V\norm{v-a}^2+\frac\mu2\norm{z-s(v)}^2$. By (a),
$\norm{v-a}^2+\norm{z-s(a)}^2\le(1+2H^2)\norm{v-a}^2+2\norm{z-s(v)}^2$, and
the choice of $g_F$ completes the proof.
\end{proof}

\begin{lemma}[Core-only tails]\label{lem:app:rec:core-tails}
Let $g_V$ be the point growth modulus of $V_\gamma$ on $[0,1]^k$ (zero if the
optimal core is not unique), and let the core coefficients be independent with
law $U_{\sigma,M}$. There is a base-computable $C_{\rm tail}$ with
$\log C_{\rm tail}\le\poly_d(I)$ such that
\begin{enumerate}[label=(\alph*)]
\item $\Pr\{g_V<t\}\le kt/\sigma+2kC_{\rm tail}/M$ for every $t>0$;
\item $\Pr\{g_V>0$ and some core coordinate at a bound at $x^*$ has
$\abs{\partial_iF_\gamma(x^*)}\le\tau\}\le N_C(\tau/\sigma+1/M)$, where
$N_C=\max\{1,k3^n\max(1,d-1)^n\}$.
\end{enumerate}
\end{lemma}

\begin{proof}
(a) The function $V_0$ is the reduced objective of $F_0$ on
$\mathcal D=[0,1]^{k+n_R}$ with the core as perturbed coordinates, and $g_V$ is
its point-growth modulus under the core noise. \ref{it:app:rec:finite} with
$p=k$, $m=n_R$, $s=2n$, and $S=k$ gives (a).

(b) A face $\Phi$ of $[0,1]^n$ fixes some coordinates at bounds and leaves a set
$\mathcal F_\Phi$ free. Let $i\in C$ be fixed on $\Phi$ and condition on all
core coefficients except $\gamma_i$. The stationary system
$\partial_jF_0(x)+\gamma_j[j\in C]=0$, $j\in\mathcal F_\Phi$, on $\Phi$ does not
involve $\gamma_i$, and its equations have degree at most $\max(1,d-1)$. By
\ref{it:app:rec:bezout} it has at most $\max(1,d-1)^n$ solutions with
nonsingular Jacobian; the empty system has the single vertex. At such a
solution $x$, $\partial_iF_\gamma(x)=\partial_iF_0(x)+\gamma_i$ lies in
$[-\tau,\tau]$ with conditional probability at most $\tau/\sigma+1/M$ by
\eqref{eq:model:interval}. A union bound over at most $3^n$ faces, $k$ core
coordinates, and the solutions gives the bound for the event that some such
solution has a small core gradient. If $g_V>0$, \cref{lem:app:rec:lift}(c)
gives full growth $g_F>0$ at $x^*$; on the face $\Phi^*$ whose relative
interior contains $x^*$, two-sided directions give
$\nabla^2_{\mathcal F\mathcal F}F_0(x^*)\succeq2g_FI$, so $x^*$ is a
nonsingular solution of the system of $\Phi^*$, and the event of (b) is
contained in the union.
\end{proof}

\begin{proof}[Proof of \cref{prop:rec:tube}]
Identify $K$ with a $q$-dimensional box by deleting its fixed coordinates.

(a) The selector $s$ is continuous by \cref{lem:app:rec:lift}(a) and
semialgebraic: by residual convexity, $y=s(v)$ is equivalent to the box
Karush--Kuhn--Tucker formula
\[
 \mathcal G(v,y):\quad v\in K\ \wedge\ \bigwedge_{i\in R}\bigl[
 (y_i=\ell_i\wedge\partial_{z_i}F_0\ge0)\vee(\ell_i\le y_i\le u_i\wedge
 \partial_{z_i}F_0=0)\vee(y_i=u_i\wedge\partial_{z_i}F_0\le0)\bigr],
\]
with derivatives evaluated at $(v,y)$. Hence each $A_i^\ell,A_i^u$ is a closed
semialgebraic subset of $K$, and $\operatorname{bd}_KA=A\setminus
\operatorname{int}_KA$ is closed and semialgebraic. It has empty interior in
$\R^q$: an open set inside it would lie in $\operatorname{int}_KA$. A
semialgebraic subset of $\R^q$ with empty interior has dimension at most $q-1$
\citep{basu2006-algorithms-in-real-algebraic-geometry}, and so does $\partial K$; hence
$\dim S_K\le q-1$, and $S_K$ is compact. The map $\nabla_KV_0$ is continuous
by \cref{lem:app:rec:lift}(b) and semialgebraic, and semialgebraic maps do not
increase dimension, so $\mathcal E_K$ is compact with $\dim\mathcal E_K\le q-1$.

Let $a\in\operatorname{relint}K$ and $0<\eta'<\dist(a,S_K)$, and let
$\mathcal B$ be the relative open ball of radius $\eta'$ about $a$. A segment
from $a$ to a point of $\mathcal B$ outside $K$ would cross $\partial K\subseteq
S_K$, so $\mathcal B\subseteq\operatorname{relint}K$. If $\mathcal B$ met both
$A_i^\ell$ and its complement, connectedness of $\mathcal B$ would give a
point of $\mathcal B$ in $\operatorname{bd}_KA_i^\ell\subseteq S_K$. So the
status of each residual coordinate is constant on $\mathcal B$.

(b) For $v\in K$, $v\in\operatorname{bd}_KA_i^\ell$ if and only if
\[
 \exists y\ \forall\varepsilon\ \exists(w,z):\quad\mathcal G(v,y)\wedge
 y_i=\ell_i\wedge\bigl[\varepsilon\le0\vee\{\mathcal G(w,z)\wedge
 \norm{w-v}^2<\varepsilon^2\wedge z_i>\ell_i\}\bigr],
\]
because $A_i^\ell$ is closed, so its relative boundary consists of its points
every relative neighborhood of which meets $K\setminus A_i^\ell$. The image
$-\nabla_KV_0(\operatorname{bd}_KA_i^\ell)$ is the set of $\beta\in\R^q$ with
\[
 \exists(v,y)\ \forall\varepsilon\ \exists(w,z):\quad
 [\text{the matrix above}]\wedge\beta+\nabla_vF_0(v,y)_K=0 .
\]
Upper boundaries use $y_i=u_i$ and $z_i<u_i$, and the image of $\partial K$
is $\exists(v,y):\mathcal G(v,y)\wedge v\in\partial K\wedge
\beta+\nabla_vF_0(v,y)_K=0$. Each of these at most $2n_R+1$ formulas has at most
three blocks of at most $N+1$ variables, $N=q+n_R\le n$, $q$ free variables, at
most $s_0=50(N+1)$ atoms, and degree at most $D_0=\max\{2,d\}$. By
\ref{it:app:rec:qe} each is equivalent to a quantifier-free formula with at
most $\mathsf Q=(s_0D_0)^{a(N+1)^3}$ disjuncts, $\mathsf Q$ atoms per
disjunct, and atom degrees at most $\mathsf Q$, for an absolute constant $a$;
this bound does not depend on coefficient sizes. Let $P_j$ be the product of
the distinct nonconstant atom polynomials of the $j$th formula ($P_j=1$ if
there are none). If an image point $\beta_0$ had $P_j(\beta_0)\ne0$, every atom
would have constant sign near $\beta_0$, the formula would hold on a
neighborhood, and the image would have nonempty interior, contradicting (a).
So each image lies in $Z(P_j)$, and $\deg P_j\le\mathsf Q^3$. The product of
the $P_j$ is a nonzero polynomial of degree at most
$D_*=(2n_R+1)\mathsf Q^3$ whose zero set contains $\mathcal E_K$, and
$\log D_*\le\poly_d(I)$. The algorithm never computes these polynomials.

(c) Let $P_K$ be the polynomial of (b). For a ball of radius $T$ and
$\epsilon\le T$, \ref{it:app:rec:tube} with $m=q-1$ gives the probability bound
$16qD_*(4D_*+2)^{q-1}\epsilon/T$ for $Z(P_K)$. The cube $[-R_0,R_0]^q$ lies in
the ball of radius $T=\sqrt qR_0$, whose volume is at most $q^{q/2}(2R_0)^q$.
Hence for $\epsilon\le R_0$ the normalized volume of
$\{x\in[-R_0,R_0]^q:\dist(x,Z(P_K))\le\epsilon\}$ is at most
$16q^{(q+1)/2}D_*(4D_*+2)^{q-1}\epsilon/R_0\le C(q,D_*)\epsilon/R_0$; since
$C\ge1$, the bound $\min\{1,C\epsilon/R_0\}$ holds for all $\epsilon\ge0$.
Now let $h=\sigma/(M-1)$ and $R_0=\sigma+h$. The cubes $g+[-h,h]^q$ centered at
the $M^q$ grid points have disjoint interiors and tile $[-R_0,R_0]^q$, and each
has volume $(2R_0)^q/M^q$. If $\dist(g,\mathcal E_K)\le\delta$, every point
of its cube is within $\delta+\sqrt qh$ of $Z(P_K)$. Therefore
\[
 \Pr\{\dist(\gamma_K,\mathcal E_K)\le\delta\}\le
 \min\Bigl\{1,C(q,D_*)\frac{\delta+\sqrt qh}{R_0}\Bigr\},\qquad
 \frac{\delta+\sqrt qh}{R_0}=\frac{M-1}M\frac\delta\sigma+\frac{\sqrt q}M .
\]
The bound includes the atoms that lie on $\mathcal E_K$.
\end{proof}

\begin{proof}[Proof of \cref{prop:rec:release}]
Let $\mathcal B$ be the relative open ball of radius $\eta$ about $a$, let
$\mathcal A$ be the residual coordinates at a bound on $\mathcal B$, with bound
values $\bar a$, and let $\mathcal F$ be the residual coordinates strictly
inside on $\mathcal B$. Write $\varsigma_i=1$ for a lower and $\varsigma_i=-1$
for an upper bound, and $v$ for the free core coordinates of $K$.

\emph{The branch.} For $v\in\mathcal B$, $s(v)=(s_{\mathcal F}(v),\bar a)$ and
$\nabla_{\mathcal F}F_0(v,s_{\mathcal F}(v),\bar a)=0$. Since
$\nabla^2_{\mathcal F\mathcal F}F_0\succeq\mu I$, the implicit function theorem
gives near each point of $\mathcal B$ a smooth solution of this equation, which
coincides with $s_{\mathcal F}$ because the strongly monotone map
$\nabla_{\mathcal F}F_0(v,\cdot,\bar a)$ has at most one zero in the box. Hence
$s_{\mathcal F}$, $V_0$, and the multipliers
$\lambda_i(v)=\varsigma_i\partial_{z_i}F_0(v,s(v))$, $i\in\mathcal A$, are smooth
on $\mathcal B$, and $\lambda_i\ge0$ there by the optimality conditions of the
residual problem.

\emph{Derivative bounds.} Let $W(v)=(v,s_{\mathcal F}(v),\bar a)$. Differentiating
$\nabla_{\mathcal F}F_0(W(v))=0$ gives
$Ds_{\mathcal F}=-(\nabla^2_{\mathcal F\mathcal F}F_0)^{-1}\nabla^2_{\mathcal Fv}F_0$,
so $\norm{Ds_{\mathcal F}}\le H$ and $\norm{DW}\le1+H$, and differentiating again
gives $\nabla^2_{\mathcal F\mathcal F}F_0\,D^2s_{\mathcal F}[p,p']
=-D^3F_0[\,\cdot_{\mathcal F},DWp,DWp']$, so
$\norm{D^2s_{\mathcal F}}\le K_3(1+H)^2/\mu$ by \eqref{eq:app:rec:hessvar}.
Since
$D^2\lambda_i[p,p']=\varsigma_i\bigl(D^3F_0[e_i,DWp,DWp']+
\nabla^2_{i\mathcal F}F_0\,D^2s_{\mathcal F}[p,p']\bigr)$,
\[
 \norm{\nabla^2\lambda_i}\le K_3(1+H)^2+K_2K_3(1+H)^2/\mu=K_3(1+H)^3\le K_\lambda
 \qquad\text{on }\mathcal B .
\]

\emph{Small multipliers have small gradients.} Let $i\in\mathcal J$, so
$0\le\lambda_i(a)\le\theta\le K_\lambda\eta^2/2$. If $\lambda_i(a)=0$, then $a$
minimizes the nonnegative function $\lambda_i$ on the open ball, so
$\nabla\lambda_i(a)=0$. Otherwise let $\beta_i=\norm{\nabla\lambda_i(a)}$, assume
$\beta_i>0$, and move from $a$ in the direction $-\nabla\lambda_i(a)/\beta_i$.
For $0<t<\eta$, Taylor's theorem and nonnegativity give
$0\le\lambda_i(a)-t\beta_i+K_\lambda t^2/2$. Letting $t$ increase to
$\min\{\eta,\sqrt{2\lambda_i(a)/K_\lambda}\}=\sqrt{2\lambda_i(a)/K_\lambda}$ gives
$\beta_i^2\le2K_\lambda\lambda_i(a)$. Hence the matrix $B$ with columns
$\varsigma_i\nabla\lambda_i(a)$, $i\in\mathcal J$, satisfies
$\norm{B}^2\le\norm{B}_F^2\le2n_RK_\lambda\theta\le\mu g/2$.

\emph{Elimination.} Near $(a,\bar a_{\mathcal J})$ define $\Pi(v,y)=F_0(v,\phi(v,y),y,
\bar a_{\mathcal A\setminus\mathcal J})$, where $\phi$ solves
$\nabla_{\mathcal F}F_0=0$ by the implicit function theorem. Its Hessian at
$(a,\bar a_{\mathcal J})$ is the Schur complement of the $\mathcal F\mathcal F$ block in the
Hessian of $F_0$ at $x^*$ restricted to $(\mathcal F,v,\mathcal J)$. Its blocks are as
follows. Since $V_0(v)=\Pi(v,\bar a_{\mathcal J})$ on $\mathcal B$, the $vv$ block is
$\nabla_K^2V_0(a)=\nabla_K^2V_\gamma(a)\succeq2gI$. By the envelope identity,
$\partial_{y_i}\Pi(v,\bar a_{\mathcal J})=\partial_{z_i}F_0(v,s(v))=\varsigma_i\lambda_i(v)$,
so the block of mixed $(v,y)$ derivatives is $B$. The $yy$ block is a Schur complement of the
residual Hessian and is $\succeq\mu I$, since for every $u$ it equals a
minimum over $w$ of the residual quadratic form at $(w,u)$. For vectors
$(p,p')$, Young's inequality gives
\[
 p^{\mathsf T}\nabla_K^2V_0p+2p^{\mathsf T}Bp'+p'^{\mathsf T}\Pi_{yy}p'
 \ge2g\norm p^2-g\norm p^2-\frac{\norm{Bp'}^2}g+\mu\norm{p'}^2
 \ge g\norm p^2+\frac\mu2\norm{p'}^2 .
\]
So $\nabla^2\Pi\succeq\delta I$ with $\delta=\min\{g,\mu/2\}$.

\emph{Restoring the eliminated coordinates.} Write the Hessian of $F_0$ at $x^*$
on $(\mathcal F,q)$, $q=(v,y)$, as $\bigl(\begin{smallmatrix}\mathsf D&\mathsf E\\
\mathsf E^{\mathsf T}&\mathsf G\end{smallmatrix}\bigr)$ with $\mathsf D\succeq\mu I$
and Schur complement $\nabla^2\Pi\succeq\delta I$. For $(w,q)$ put
$\hat w=w+\mathsf D^{-1}\mathsf Eq$; then the quadratic form equals
$\hat w^{\mathsf T}\mathsf D\hat w+q^{\mathsf T}\nabla^2\Pi\,q\ge
\mu\norm{\hat w}^2+\delta\norm q^2$. Since
$\norm{\mathsf D^{-1}\mathsf E}\le K_2/\mu=H$,
$\norm w^2+\norm q^2\le2\norm{\hat w}^2+(1+2H^2)\norm q^2$, and therefore the
form is at least $\nu(\norm w^2+\norm q^2)$. Principal submatrices of a matrix
$\succeq\nu I$ are $\succeq\nu I$.
\end{proof}

\begin{proof}[Proof of \cref{thm:rec:core-only}]
\emph{Base choices.} Let $B$ be the factor of \ref{it:app:rec:fallback} for
$F_0$ on $[0,1]^n$ (it applies to every rational linear term, in particular to
core-only noise), enlarged to bound $q$-bit evaluation. Let $C_{\rm tail}$ and
$N_C$ be as in \cref{lem:app:rec:core-tails}, let $D_*$ be as in
\cref{prop:rec:tube}, and put $C_S=3^k\max_{1\le q\le k}C(q,D_*)$. Choose
\begin{equation}\label{eq:app:rec:core-params}
\begin{gathered}
 \rho=\frac1{6B},\quad g=\frac{\rho\sigma}{2k},\quad
 \tau=\frac{\rho\sigma}{2N_C},\quad
 \eta=\min\Bigl\{\frac14,\frac{\rho\sigma}{4C_SH_V}\Bigr\},\quad
 K_\lambda=\max\{1,K_3(1+H)^3\},\\
 \theta=\min\Bigl\{\frac{K_\lambda\eta^2}2,\frac{\mu g}{4n_RK_\lambda}\Bigr\},\quad
 \nu=\min\Bigl\{\frac\mu2,\frac{\min\{g,\mu/2\}}{1+2H^2}\Bigr\},\quad
 g_*=\min\Bigl\{\frac g{1+2H^2},\frac\mu4,\frac\nu4\Bigr\},\\
 r=\min\Bigl\{\frac18,\frac\tau{16K_2},\frac\theta{16K_2},\frac{g_*}{4K_3}\Bigr\},\quad
 A_0=2+\frac{kL+8}{g_*},
\end{gathered}
\end{equation}
let $J$ be the least $j\ge0$ with
$2^{-j}\le\min\{r/(4A_0),g_*r^2/(16GA_0)\}$, and let $M$ be the least power of
two with $M\ge\max\{2,2^J,4kC_{\rm tail}/\rho,2N_C/\rho,2kC_S/\rho\}$. All these
are base-computable with polynomial bit length, and $J,\log_2M\le\poly_d(I)$.

\emph{Algorithm and correctness.} Run the algorithm of \cref{app:rec:main} in
certified mode with $g_*$ in place of $g_0$ and the parameters above: the core
search with accuracy $E_j$, slab bounds with accuracy $g_*r^2/16$, sign tests on
all original bounds (core and residual), and the Hessian test with modulus
$g_*$. Every residual coordinate remains eligible for fixing only at an original
bound. The fallback of \ref{it:app:rec:fallback} runs on the same draw if no
level returns. Correctness on every draw follows exactly as in
\cref{app:rec:main}.

\emph{The good event.} Let $\mathcal B_1=\{g_V<g\}$, let $\mathcal B_2$ be the
event of \cref{lem:app:rec:core-tails}(b), and let $\mathcal B_3$ be the event
that $g_V>0$ and the optimal core $a$ lies in the relative interior of a face
$K$ with $q\ge1$ free coordinates and $\dist(a,S_K)\le2\eta$. On $\mathcal B_3$,
$a$ minimizes $V_\gamma$ along $K$ and $V_\gamma$ is differentiable, so
$\gamma_K=-\nabla_KV_0(a)$; for a nearest point $\hat v\in S_K$,
$\dist(\gamma_K,\mathcal E_K)\le\norm{\nabla_KV_0(a)-\nabla_KV_0(\hat v)}
\le2H_V\eta$. By \cref{lem:app:rec:core-tails}, \cref{prop:rec:tube}(c), a union
over the at most $3^k$ faces of the core box, and the choices above,
\[
 \Pr(\mathcal B_1)\le\frac{kg}\sigma+\frac{2kC_{\rm tail}}M\le\rho,\quad
 \Pr(\mathcal B_2)\le N_C\Bigl(\frac\tau\sigma+\frac1M\Bigr)\le\rho,\quad
 \Pr(\mathcal B_3)\le C_S\Bigl(\frac{2H_V\eta}\sigma+\frac kM\Bigr)\le\rho .
\]
So the bad event has probability at most $3\rho=1/(2B)$.

\emph{Stopping on the good event.} Outside $\mathcal B_1\cup\mathcal B_2\cup
\mathcal B_3$ the optimal core $a$ is unique with projected growth at least
$g$, and by \cref{lem:app:rec:lift}(c) $x^*=(a,s(a))$ has full point growth at
least $g_F\ge g_*$. The hull, incumbent, and containment estimates of
\cref{lem:app:rec:stopping} hold with $g_*$ in place of $g_0$ (the hull estimate
uses the projected growth $g\ge g_*$), so at level $J$ step~(2) passes and every
point of $Q$ is within $5r/4$ of $x^*$. The sign tests fix every core coordinate
at a bound, since its derivative exceeds $\tau$ in magnitude, and every active
residual coordinate with multiplier $\lambda_i>\theta$, since the midpoint test
value is at least $\lambda_i-\frac94K_2r>\theta-\frac9{64}\theta$. Coordinates
strictly inside their bounds at $x^*$ are not fixed. If $q\ge1$, outside
$\mathcal B_3$ the residual status is constant on the relative ball of radius
$\eta$ about $a$, $\nabla_KV_\gamma(a)=0$, and growth gives
$\nabla^2_KV_\gamma(a)\succeq2gI$; \cref{prop:rec:release} then gives a Hessian at
least $\nu I$ on the unfixed coordinates. If $q=0$, all core coordinates are
fixed and the unfixed coordinates are residual, with Hessian at least $\mu I\succeq
\nu I$. In both cases $\nu\ge4g_*$, and \eqref{eq:app:rec:hessvar} gives
$\nabla^2F_0(m')-(K_3r_{Q'}+g_*)I\succeq(4g_*-2K_3r-g_*)I\succeq\frac52g_*I$.
The algorithm returns by level $J$.

\emph{Work and outputs.} The expected-count bound of \cref{prop:rec:search}
needs only the independent core coefficients and $M\ge2^J$. The rest of the
accounting, the proof record, and the evaluation bounds are as in the proof of
\cref{thm:rec:poly}.
\end{proof}

\subsection{Examples}\label{app:rec:examples}

\begin{proof}[Verification of \cref{ex:rec:feasibility}]
For each $v$, the least feasible $z$ is $\abs{v-\frac13}\le\frac23$, so
$V(v)=\abs{v-\frac13}$. Since $2^j\not\equiv0\pmod3$, the number $2^j/3$ has
fractional part $\frac13$ or $\frac23$, so $\frac13$ lies in the interior of a
level-$j$ cell whose endpoints are at distances $h_j/3$ and $2h_j/3$ from it.
Every grid value is at least $h_j/3$, so $U_j-E_j\ge h_j/3-kLh_j^2/8>0=f^*$ when
$h_j<8/(3kL)$.
\end{proof}

\begin{proof}[Verification of \cref{ex:rec:star} and its variant]
Each leaf term $z^2-\frac v2z$ is minimized over $[0,1]$ at $z=v/4$ with value
$-v^2/16$, which gives $V$. Over $\{0,\frac12,1\}$ its minimum is $0$ for
$v\in\{0,\frac12,1\}$ (attained at $z=0$, and also at $z=\frac12$ when $v=1$),
so $V_h(v)=v^2+\frac{15}{16}v$ on the grid, which gives the three values, and
$V_h(v)-V(v)=m\min_{z}(z-v/4)^2$. The four bag corners
$(v,z_1)\in\{\frac12,1\}\times\{0,\frac12\}$ have min-marginals
$\frac{23}{32}$, $\frac{23}{32}+\frac18$, $\frac{31}{16}$, $\frac{31}{16}$, and
the grid minimum is $V_h(0)=0$. Since $\frac{23}{32}-\frac18>0$, the bag-local
rule discards the cell, while $\frac{23}{32}-\frac{33}{16}<0$.

For the variant let $d\ge8$, $m=d^2$, $h=\frac14$, and
$F_0(v,z)=(v-\frac12)^2+\sum_{i=1}^m(z_i-\frac18-(v-\frac12)/d)^2$ on
$[0,1]^{m+1}$. Its Hessian has eigenvalue $2$ with multiplicity $m-1$ and
eigenvalues $3\pm\sqrt5$ on the span of $e_v$ and the normalized sum of the leaf
directions, so it is uniformly positive definite, $L=4$ is valid, and the star
graph has coloring number two. With $t=v-\frac12$, every ideal leaf value
$\frac18+t/d$ lies in $[\frac1{16},\frac3{16}]$, its nearest point of
$\{0,\frac14,\ldots,1\}$ is $0$ or $\frac14$ at distance $\frac18-\abs t/d$, and
\[
 V(v)=t^2,\qquad V_h(v)=\frac{d^2}{64}-\frac{d\abs t}4+2t^2 .
\]
The grid minimum is $U=\frac{d^2}{64}-\frac d8+\frac12$ at $t=\pm\frac12$. The
cell $v\in[\frac12,\frac34]$, $z_1\in[0,\frac14]$ contains the optimizer
$(\frac12,\frac18,\ldots,\frac18)$, and its least corner min-marginal is
$q_C=\frac{d^2}{64}-\frac d{16}+\frac18$, at $(v,z_1)=(\frac34,\frac14)$. Then
$q_C-U=\frac d{16}-\frac38\ge\frac18>\frac1{16}=\abs BLh^2/8$, so the bag-local
rule discards the cell. Finally $(V_h-V)(\frac34)-(V_h-V)(\frac12)=-(d-1)h^2$.
\end{proof}

\begin{proof}[Verification of \cref{ex:rec:fiber}]
Write $\varrho=z_1-vz_2$. The residual Hessian is
$2(1,-v)^{\mathsf T}(1,-v)\succeq0$ and $\partial_{vv}F_0=2+2z_2^2\le4$. Taking
$z=(0,0)$ shows $V_\gamma(v)=(v-\frac12)^2+\gamma v=(v-a_\gamma)^2+f^*$, which
gives the optimal core, its growth constant, and the optimal set
$\{\varrho=0,v=a_\gamma\}$. For $p=(0,v,1)$ one computes
$p^{\mathsf T}\nabla^2F_0\,p=0$ and $\nabla^2F_0\,p=(-2\varrho,0,0)$. A positive
semidefinite matrix with $p^{\mathsf T}Hp=0$ has $Hp=0$, so the Hessian is not
positive semidefinite where $\varrho\ne0$; every full-dimensional box contains
such points.

For the strictly convex variant put $\alpha=z_1-\frac12$, $\beta=z_2-\frac12$, and
$\varrho=\alpha-v\beta$. The residual problem $\min(\alpha-v\beta)^2+\beta^4$ is
strictly convex: on a segment either $\beta$ varies, and $\beta^4$ is strictly
convex, or $\beta$ is constant, and the square is strictly convex in $\alpha$. Its
unique minimizer is $\alpha=\beta=0$, so $V_\gamma$ and $a_\gamma$ are unchanged and
the optimizer $(a_\gamma,\frac12,\frac12)$ is unique and interior. At points
with $\beta=0$ and $\alpha\ne0$, the vector $p=(0,v,1)$ in the coordinates
$(v,\alpha,\beta)$ satisfies $p^{\mathsf T}Hp=0$ and $Hp=(-2\alpha,0,0)$, and the
direction $p+\alpha e_v$ has $(p+\alpha e_v)^{\mathsf T}H(p+\alpha e_v)=-2\alpha^2<0$.
Such points lie arbitrarily close to the optimizer.
\end{proof}

\begin{proof}[Verification of \cref{ex:rec:weak,ex:rec:two-sided}]
For $F_0=\phi(v)+y^2$, $\partial_yF_0=2y\ge0$ with equality only at $y=0$, so the
selector is $0$ and the multiplier $\partial_yF_0(v,0)$ vanishes for all $v$.
For the coupled variant, $\partial_yF_0=2(y+(v-\frac12)^2)\ge0$ on $y\ge0$, the
selector is $0$, the multiplier is $2(v-\frac12)^2$, and
$V_0(v)=(v-\frac12)^2+(v-\frac12)^4$.

In \cref{ex:rec:two-sided}, $\partial_yF_\gamma=v+y\ge0$ on the box, with
equality only at $(0,0)$, so the selector is $0$ and
$V_\gamma(v)=\frac14v^2+\gamma v$. The Hessian has determinant $\frac12-1<0$. The
multiplier is $\lambda(v)=v$. For $\gamma>0$, $V_\gamma$ is increasing on $[0,1]$,
so the optimum is $(0,0)$ with $\partial_vF_\gamma(0,0)=\gamma$.
\end{proof}
